\section{Appendix to Section \ref{sec:syntactic}}\label{app:sec:syntactic}


%%%%%%%%COMMENTO LA DIMOSTRAZIONE FI TC CBC PERCHè ORA SEMPLIFICATA E NEL TESTO%%%%%%%%%%%%%
\begin{comment}
    
\begin{proof}[Proof of Theorem \ref{thm:TCconvexbiproductcategory}]
The fact that the definitions of $+_p$ and $\star$ in \eqref{eq:enrichmentTC} provides a  $\Cat{PCA}$ follows immediately by the axioms \ref{eq:diagp assoc}, \ref{eq:diagp idempotency} and \ref{eq:diagp symmetry}. The fact that this provides an enrichment of $\CatTapeC$ over $\Cat{PCA}$ is easily proved by using naturality of $\diagp{}$, $\codiag{}$, $\bangp{}$ and $\cobang{}$ see, e.g., \cite{bonchi2025tapediagramsmonoidalmonads}.

The fact that $\zero$ is both initial and final follows immediately by naturality of  $\cobang{}$ and $\bang{}$: axioms \eqref{eq:bangp nat} and \eqref{eq:cobang nat}.

Since, by the axioms in Figure \ref{fig:freestrictfccat}, every object $P$ of $\CatTapeC$ is equipped with a natural and coherent monoid, by \cite{fox1976coalgebras}, it holds that $\oplus$ is a coproducts with injections
\begin{equation*}
\iota_1 \defeq \id{P_1}\oplus \cobang{P_2} \colon P_1 \to P_1\oplus P_2\quad\text{and}\quad\iota_2 \defeq \cobang{P_1}\oplus \id{P_2} \colon P_2 \to P_1\oplus P_2\text{.}
\end{equation*}
Recall that (see for instance \cite[Ch. 6.4]{mellies2009categorical} ) the copairing of $\t\colon P_1 \to Q$ and $\s\colon P_2 \to Q$ is defined as
\begin{equation*}
[\t,\s] \defeq (f\oplus g); \codiag{Q} 
\end{equation*}
and thus $[\id{Q},\id{Q}] = \codiag{Q}$.

Our main technical effort, illustrated below in Proposition \ref{prop:TChasconvexcoprduct}, consists in proving that $\oplus$ is also a convex product with projections
\begin{equation*}
\pi_1 \defeq \id{P_1}\oplus \bang{P_2} \colon P_1 \oplus P_2 \to P_1\quad\text{and}\quad\pi_2 \defeq \bang{P_1}\oplus \id{P_2} \colon P_1\oplus P_2 \to  P_2\text{.}
\end{equation*}
and the convex pairing of $\t\colon P \to Q_1$ and $\s\colon P \to Q_2$ is given as
\begin{equation*}
<\t,\s>_{1,0}\defeq  \t\oplus \cobang{Q_2} \qquad <\t,\s>_{0,1} \defeq \cobang{Q_1} \oplus \s \qquad <\t,\s>_{p,1-p} \defeq \; \diagp{P}; (\t\oplus \s) \text{ for }p\in(0,1)
\end{equation*}
As we will illustrate later, the case of  $<\t,\s>_{p,q}$ with $p+q<1$ simply follows from the above one.

Before delving in proof of Proposition~\ref{prop:TChasconvexcoprduct} observe that \eqref{eq:delta} in Definition~\ref{def:convbicat} follows easily by the above definitions of $\iota_i$ and $\pi_1$, axioms~\eqref{eq:bangp nat} and \eqref{eq:cobang nat}, for the case $i=j$, and the definition of $\star$ in \eqref{eq:enrichmentTC} for the case $i\neq j$.
%Instead the laws \eqref{eq: assioma aggiuntivo} follows from the characterizations of $<f,g>_{p,1-p}$ and $\codiag{}$ and the definition of $+_p$ in \eqref{eq:enrichmentTC}.
Finally, to prove that $\langle \t_1, \dots \t_n\rangle_{\vec{p}} = \diagpn{\;\;\vec{p}}{U}{n}; \bigoplus_{i=1}^n \t_i$, one can easily check that, for all $i$, $(\diagpn{\;\;\vec{p}}{U}{n}; \bigoplus_{i=1}^n \t_i ); \pi_i =p_i\cdot \t_i$, where the $i$-th projection $\pi_i\colon \bigoplus_{j=1}^nU_j \to U_i$ is defined as $(\bang{\bigoplus_{j=1}^{i-1}U_j}) \oplus \id{U_i} \oplus (\bang{\bigoplus_{j=i+1}^{n}U_j})$.
\end{proof}

The proof of Proposition~\ref{prop:TChasconvexcoprduct} substantially relies on a normal form result (Lemma~\ref{lemma:division}) and on the fact that the enrichment on $\CatTapeC$ is cancellative, namely if $p \cdot \t = p \cdot \s$, then $\t=\s$ (Lemma~\ref{cancellativity}). The remainder of this appendix is structured as follows: we first illustrate some preliminary results, then some results about normal forms, then the cancellativity property and finally convex products.
\end{comment}


\subsection{Preliminary Results}
We collect below several simple results that we will frequently use.

\begin{lemma}\label{lemma:pcdot}
Let $\t\colon P \to Q$ be an arrow of $\CatTapeC$ . Then, for all $p\in(0,1)$, $p\cdot \t =\, \diagp{P};(\t \oplus \bang{P})$.
\end{lemma}
\begin{proof}
\begin{align*}
p\cdot \t &= \t+_p \star_{P,Q} \tag{def. of $p\cdot -$}\\
&=\, \diagp{P} ; (\t \oplus \star_{P,Q}) ; \codiag{Q} \tag{\ref{eq:enrichment}}\\
&=\, \diagp{P} ; (\t \oplus (\bang{P}; \cobang{Q})); \codiag{Q} \tag{\ref{eq:enrichment}}\\
&=\, \diagp{P} ; (\t \oplus  \bang{P}) ;(\id{Q} \oplus \cobang{Q}) ; \codiag{Q} \tag{Symmetric Monoidal Category}\\
&=\, \diagp{P} ; (\t \oplus  \bang{P}) ;\id{Q} \tag{\ref{eq:codiag unital}} \\
&=\, \diagp{P} ; (\t \oplus  \bang{P})  \tag{Category} 
\end{align*}
\end{proof}


\begin{lemma}\label{lemma:pcopairing}
Let $\s = (\s_1 \oplus \s_2) ; \codiag{}$ be an arrow of $\CatTapeC$ . 
Then $r \cdot \s = (r\cdot \s_1 \oplus r\cdot \s_2); \codiag{}$. 
\end{lemma}
\begin{proof}
Assume that $\s$ has type $P_1\oplus P_2 \to Q$ and observe that
\begin{align*}
r \cdot \s &=\, \diagpX{r}{P_1\oplus P_2} ; (\s \oplus \bang{}) \tag{Lemma \ref{lemma:pcdot}}\\
&= (\diagpX{r}{P_1} \oplus\, \diagpX{r}{P_2} ) ; (\id{P_1}\oplus \symm{P_1}{P_2} \oplus \id{P_2}) ; (\s \oplus \bang{P_1}\oplus \bang{P_2}) \tag{\ref{eq:diagp coherence}}\\
&= (\diagpX{r}{P_1} \oplus\, \diagpX{r}{ P_2} ) ; (\id{P_1}\oplus \symm{P_1}{P_2} \oplus \id{P_2}) ; (\, ((\s_1 \oplus \s_2) ; \codiag{Q}) \oplus \bang{P_1}\oplus \bang{P_2}\,) \tag{Hypothesis}\\
&= (\diagpX{r}{P_1} \oplus\, \diagpX{r}{P_2} ) ; (\id{P_1}\oplus \symm{P_1}{P_2} \oplus \id{P_2}) ; ( \s_1 \oplus \s_2  \oplus \bang{P_1}\oplus \bang{P_2}); \codiag{Q} \tag{Sym. Mon. Cat.}\\
&= (\diagpX{r}{P_1} \oplus\, \diagpX{r}{ P_2} ) ;( \s_1 \oplus \bang{P_1} \oplus \s_2   \oplus \bang{P_2}); \codiag{Q} \tag{Symmetric Monoidal Category}\\
&= (\diagpX{r}{P_1} ; ( \s_1 \oplus \bang{P_1})) \oplus (\diagpX{r}{ P_2}  ;( \s_2   \oplus \bang{P_2})); \codiag{Q} \tag{Symmetric Monoidal Category}\\
&= (r\cdot \s_1 ) \oplus (r\cdot  \s_2 ); \codiag{Q} \tag{Lemma \ref{lemma:pcdot}}
\end{align*}
\end{proof}


\begin{lemma}\label{lemma:basicTC}
Let $\t=\,\diagp{} ; (\t_1 \oplus \t_2)$ be an arrow of $\CatTapeC$. Then:
\begin{enumerate}
\item $\t; (\id{} \oplus \bang{}) = p\cdot \t_1$;
\item $\t ; (\bang{} \oplus \id{}) = (1-p) \cdot \t_2$;
\item $r \cdot \t =\, \diagp{} ; (r\cdot \t_1 \oplus r\cdot \t_2)$ for all $r\in (0,1)$.
\end{enumerate}
\end{lemma}
\begin{proof}
Points 1 and 2 follow immediately from the definition of the enrichment.

We prove below point 3 for $\t_1\colon:P\to Q_1$ and $\t_2\colon P\to Q_2$.
\begin{align*}
r \cdot \t &=\, \diagpX{r}{} ; (\t \oplus \bang{Q_1\piu Q_2}) \tag{Lemma \ref{lemma:pcdot}}\\
&=\, \diagpX{r}{} ; (\diagp{} ; (\t_1 \oplus \t_2) \oplus \bang{{Q_1\piu Q_2}}) \tag{Hypothesis}\\
&=\, \diagpX{r}{} ; ( \, (\diagp{} ; (\t_1 \oplus \t_2)) \oplus (\diagp{}; (\bang{Q_1} \oplus \bang{Q_2}))\, ) \tag{\ref{eq:bangp coherence} + \ref{eq:diagp idempotency}}\\
&=\, \diagpX{r}{} ; (\diagp{} \oplus \diagp{}) ; (\,  (\t_1 \oplus \t_2) \oplus (\bang{Q_1} \oplus \bang{Q_2}) \, ) \tag{Symmetric Monoidal Category}\\
&=\, \diagp{} ; (\diagpX{r}{} \oplus \diagpX{r}{}) ; (\id{} \oplus \symm \oplus \id{}) ; (\,  (\t_1 \oplus \t_2) \oplus (\bang{Q_1} \oplus \bang{Q_2}) \, ) \tag{\ref{eq:diagp nat} + Sym. Mon. Cat.}\\
&=\, \diagp{} ; (\diagpX{r}{} \oplus \diagpX{r}{}) ;  (\,  (\t_1 \oplus \bang{Q_1}) \oplus (\t_2 \oplus \bang{Q_2}) \, ) \tag{Symmetric Monoidal Category}\\
&=\, \diagp{} ; ( \,(\diagpX{r}{} ; (\t_1 \oplus \bang{Q_1})) \oplus (\diagpX{r}{} ;(\t_2 \oplus \bang{Q_2})) \, ) \tag{Symmetric Monoidal Category}\\
&=\, \diagp{} ; (r\cdot \t_1 \oplus r\cdot \t_2)  \tag{Lemma \ref{lemma:pcdot}}
\end{align*}
 
\end{proof}

\begin{lemma}\label{lemma:fractions}
Let $\t=\,\diagp{} ; (\t_1 \oplus \t_2)$ and $\s=\, \diagq{} ; (\s_1 \oplus \s_2)$. If $\t_1=\frac{q}{p}\cdot \s_1$ and $\frac{1-p}{1-q} \cdot \t_2 = \s_2$, then $\t = \s$.
\end{lemma}
\begin{proof}
\begin{align}
\t &=\,  \diagp{};( \frac{q}{p} \cdot \s_1 \oplus \t_2  )\tag{Lemma \ref{lemma:basicTC}.1}\\
&=\, \diagp{}; (\diagpX{\frac{q}{p}\,\,}{} \oplus \id{}) ; (\s_1 \oplus \bang{} \oplus \t_2) \tag{Lemma \ref{lemma:pcdot}}\\
&=\, \diagq; (\id{} \oplus\, \diagpX{\frac{p-q}{1-q}\,\,}{}) ; (\s_1 \oplus \bang{} \oplus \t_2) \tag{\ref{eq:diagp assoc}} \\
&=\, \diagq; (\id{} \oplus ( \diagpX{\frac{1-p}{1-q}\,\,}{}; \symm)) ; (\s_1 \oplus (\bang{} \oplus \t_2)) \tag{\ref{eq:diagp symmetry}} \\
&=\, \diagq; (\id{} \oplus ( \diagpX{\frac{1-p}{1-q}\,\,}{})) ; (\s_1 \oplus  \t_2 \oplus \bang{}) \tag{Symmetric Monoidal Category}\\
&=\,  \diagq{};( \s_1 \oplus \frac{1-p}{1-q} \cdot \t_2  ) \tag{Lemma \ref{lemma:basicTC}.2}\\
&=\s \tag{Hypothesis}
\end{align}
\end{proof}


\subsection{Normal forms}\label{app:normal forms}
We are now going to prove several normal form results. For these results we need to fix some notation.



%For all natural numbers $n \in \mathbb{N}$ and $p_1, \dots , p_n \in (0,1)$ such that $\sum_{i=1}^n p_i \leq 1$ one can define $\Br{+_{\vec{p}}}\colon A \to \bigoplus_{i=1}^n  A$ as follows:\\
%for $n=0$ take $\Br{{+_{\vec{p}}}}\coloneqq\Br{\star}_A$ ;\\
%for $n+1$ take $\Br{{+_{\vec{p}}}}\coloneqq \Br{+_{p_{n+1}}}_A; (\id{A} \oplus \Br{{+_{\frac{\vec{p'}}{1-p_{n+1}}}}})$ where $\vec{p'} = p_1, \dots, p_n$.

%\marginpar{per ottenere che la probabilità del primo ingresso sia $p_1$ e così via ed evitare di avere la $\star$ nella definizione si deve riscrivere in questo modo}

%For all natural numbers $n \in \mathbb{N}$ and $p_1, \dots , p_n \in (0,1)$ such that $\sum_{i=1}^n p_i \leq 1$ one can define $\diaggen{\vec{p}}{P}\colon P \to \bigoplus_{i=1}^{n+1}  P$ as follows:\\
%\begin{equation}\label{eq:diag vettore in appendice}
%    (n=0):\  \diaggen{\vec{p}}{P}\defeq\id{U}\qquad
%(n=k+1):\ \diaggen{\vec{p}}{P}\defeq\, \diaggen{p_1}{P}; (\id{P} \oplus\, \diaggen{\frac{\vec{p'}}{1-p_1}}{P})
%\end{equation}
%where $\vec{p'} = p_2, \dots, p_n$.
%\newline
%Observe that $\diaggen{\vec{p}}{P}$ differs from the definition of $\diagpn{\vec{p}\;}{P}{n}$ given in \eqref{eq: diagpn cap equational presentation}, but we need this formulation for the rest of the section. 
%
%
For all natural numbers $n \in \mathbb{N}$ one can define $\codiag{U}^n\coloneqq \bigoplus_{i=1}^n U\to U$ as follows:
\begin{equation}\label{eq:codiagn in appendice}
    (n=0):\  \codiag{U}^0\defeq\cobang{U}\qquad (n+1):\ \codiag{U}^{n+1}\defeq(\id{U} \oplus \codiag{U}^n  );\codiag{U}.
\end{equation}







\begin{lemma}[Pre normal form]\label{lemma:prenormalform}%\marginpar{Questa poi avra' un analogo nelle matrici}
Every arrow $\t\colon \oplus_{i=1}^n U_i\to \oplus_{i=1}^m V_i$ in $\CatTapeC$ can be written as $\t_1;\t_2;\t_3;\t_4$ where:
\begin{itemize}
    \item $\t_1$ is obtained through the fragment \setlength{\tabcolsep}{4pt}\begin{tabular}{rc ccccccccccccccccccccc}
   $\!\!\! \mid \!\!\!$ & $\diagp{U}$  & $\!\!\! \mid \!\!\!$  & $\id{\zero}$ & $\!\!\! \mid \!\!\!$ &  $\id{U}$ & $\!\!\! \mid \!\!\!$   & $   \t ; \t   $ & $\!\!\! \mid \!\!\!$  & $  \t \piu \t  $ & $\!\!\! \mid \!\!\!$  &  $\sigma_{U,V}^{\piu}$
\end{tabular}
\item $\t_2$ is obtained through the fragment \setlength{\tabcolsep}{4pt}\begin{tabular}{rc ccccccccccccccccccccc}
           $\!\!\! \mid \!\!\!$  &  $\bangp{U}$& $\!\!\! \mid \!\!\!$&  $\id{\zero}$ & $\!\!\! \mid \!\!\!$ & $\id{U}$    & $\!\!\! \mid \!\!\!$   & $   \t ; \t   $ & $\!\!\! \mid \!\!\!$  & $  \t \piu \t  $ & $\!\!\! \mid \!\!\!$  &  $\sigma_{U,V}^{\piu}$
    \end{tabular}
    \item $\t_3$ is obtained through the fragment \begin{tabular}{rc ccccccccccccccccccccc}\setlength{\tabcolsep}{0.0pt}
        $\!\!\! \mid \!\!\!$ & $\tapeFunct{c}$  & $\!\!\! \mid \!\!\!$  &  $\id{\zero}$ & $\!\!\! \mid \!\!\!$ & $\id{U}$&  $\!\!\! \mid \!\!\!$   & $   \t ; \t   $ & $\!\!\! \mid \!\!\!$  & $  \t \piu \t  $ & $\!\!\! \mid \!\!\!$  &  $\sigma_{U,V}^{\piu}$
    \end{tabular}
    \item $\t_4$ is obtained through the fragment \begin{tabular}{rc ccccccccccccccccccccc}%\setlength{\tabcolsep}{4pt}
        $\!\!\! \mid \!\!\!$ & $\codiag{U}$ &$\!\!\! \mid \!\!\!$ & $\cobang{U}$ & $\!\!\! \mid \!\!\!$  &  $\id{\zero}$ & $\!\!\! \mid \!\!\!$ & $\id{U}$ & $\!\!\! \mid \!\!\!$   & $   \t ; \t   $ & $\!\!\! \mid \!\!\!$  & $  \t \piu \t  $ & $\!\!\! \mid \!\!\!$  &  $\sigma_{U,V}^{\piu}$
    \end{tabular}
\end{itemize}
\end{lemma}
\begin{proof}
Using the naturality, one can first move all the occurrences of $\diagp{}$ to the left. Then move all the occurrences of $\cobang{}$ to the left until an occurrence of $\diagp{}$ is met (note that since $\cobang{}$ is not the unit of $\diagp{}$, $\cobang{}$ remains at the right of $\diagp{}$). Similarly, one can move all occurrences of $\codiag{}$ and $\cobang{}$ on the right.
\end{proof}




\begin{lemma}\label{lemmaABform}
Let $\t\colon U \to V$ be an arrow of $\CatTapeC$. Then one of the following holds: 
\begin{itemize}
\item $\t=\star_{U,V}$;
\item there exist  $c \in \Cat{C}[U,V]$ such that $\t=\tapeFunct{c}$;
%\item there exist distinct $c_1, \dots, c_n \in \Cat{C}[U,V]$ and $\vec{p}=p_1, \dots, p_{n+1}$ such that
%\[\t=\, \diagvecp{A} ; ( (\bigoplus_{i=1}^n  \tapeFunct{c_i}) \oplus\, \bangp{U}  ) ; \codiag{V}^n\text{.}\]
\item there exist distinct $c_1, \dots, c_n \in \Cat{C}[U,V]$ and $\vec{p}=p_1, \dots, p_{n}\in (0,1)$ such that
\[\t=\, \diagpn{\vec{p}\;}{U}{n}
;  (\bigoplus_{i=1}^n  \tapeFunct{c_i})  ; \codiag{V}^n\text{.}\]
\end{itemize}
\end{lemma}
\begin{proof}
By Lemma~\ref{lemma:prenormalform}, $\t= \t_1;\t_2;\t_3;\t_4$.
Using the naturality of $\sigma_{U,V}^{\piu}$ (see Table~\ref{fig:freestricmmoncatax}), one can move all the occurrences of $\sigma_{U,V}^{\piu}$ to the right until an occurrence of $\codiag{V}$ is met. 
Using the associativity and the symmetry of $\diagp{U}$,  all multiple occurrences of every $c_i$ and $\star_{U,V}$ can be collected in the form $\diagp{U};(c_i\otimes c_i)$ and then collapsed to one occurrence of $c_i$ by using the idempotency axiom and axioms for $\codiag{U}$. Then again using the associativity and the symmetry of $\diagp{U}$ one can move all the occurrences of $\diagp{U}$ in the form of $\diagpn{\vec{p}\;}{U}{n}$ for some vector $\vec{p}=(p_1,\dots,p_n)$, and similarly all occurrences of $\codiag{V}$ in the form of $\codiag{V}^n$.
\end{proof}

\begin{lemma}\label{lemma: arrows A-B subdistributions}
    Arrows in $\CatTapeC$ of the form $\t\colon U\to V$ are in bijection with $\mathcal{D}_{\le}(\Cat{C}[U,V])$.
\end{lemma}
\begin{proof}
    By Lemma~\ref{lemmaABform} every arrow $\t\colon U\to V$ is in one of the following forms:
\begin{itemize}
\item $\t=\star_{U,V}$;
\item there exists $c \in \Cat{C}[U,V]$ such that $\t=\tapeFunct{c}$;
\item there exist distinct $c_1, \dots, c_n \in \Cat{C}[U,V]$ and $\vec{p}=p_1, \dots, p_{n}\in(0,1)$ such that 
\[\t=\,\diagpn{\vec{p}\;}{U}{n}
;  (\bigoplus_{i=1}^n  \tapeFunct{c_i})  ; \codiag{V}^n\text{.}\]
\end{itemize}
The first two cases correspond respectively to the zero subdistribution and the Dirac distribution on $c$. The third case corresponds to the subdistribution $\vec{p}$ on the set of arrows $\{c_1, \dots, c_n\}$, where $p_i$ is the probability of $c_i$. %and $p_{n+1}=1-\sum_{i=1}^{n}p_i$ is the probability of $\Br{\star}_A$.
 If $\t$ and $\t'=\, \diagpn{\vec{q}\;}{U}{n}; (\bigoplus_{i=1}^{n'}  \tapeFunct{c'_i}) ; \codiag{V}^{n'}$ are sent to the same subdistribution, then one has that $n=n'$, $p_i=q_i$ and $c_i=c'_i$ for all $i=1,\dots,n$. 
\end{proof}

\begin{lemma}\label{lemma:division}
Let $\t \colon U \to \bigoplus_{i=1}^n V_i$ be an arrow of $\CatTapeC$. Then, 
\begin{itemize}
\item either $\t= \bang{U}; \cobang{\bigoplus_{i=1}^n V_i}$,
\item or there exists $j\in \{1,\dots, n\}$ and $\t_j \colon U \to V_j$ such that $\t = \cobang{\bigoplus_{i=1}^{j-1} V_i} \oplus \t_j \oplus \cobang{\bigoplus_{i=j+1}^{n} V_i}$,
\item or for all $j\in \{1,\dots, n\}$, there exists $\s_1\colon U \to \bigoplus_{i=1}^{j} V_i$ and $\s_2 \colon U \to \bigoplus_{i=j+1}^{n} V_i$ such that $\t=\,\diagp; (\s_1\oplus \s_2)$.
\end{itemize}
\end{lemma}
\begin{proof}
We have that $\t= \t_1;\t_2;\t_3;\t_4$ where the $\t_i$ are prescribed by Lemma~\ref{lemma:prenormalform}.
We consider two different cases: either $\t_1=\id{U}$ or $\t_1 \neq \id{U}$.

Assume that $\t_1=\id{U}\colon U \to U$. Then $\t_2$ is an arrow with domain $U$ and thus either $\t_2 = \bang{U}$ or $\t_2= \id{U}$. 

\begin{itemize}
\item If $\t_2 = \bang{U} \colon U \to \zero$, then $\t_3$ is an arrow with domain $\zero$. This entails that $\t_3$ should be $\id{\zero}\colon \zero \to \zero$ and $\t_4$ an arrow of type $\zero \to \bigoplus_{i=1}^n V_i$. The latter fact entails that $\t_4 = \cobang{\bigoplus_{i=1}^n V_i}$. Thus $\t= \bang{U}; \cobang{\bigoplus_{i=1}^n V_i}$.
\item If $\t_2 = \id{U} \colon U \to U$, then $\t_3$ should have type $U\to V_j$ and thus it must be either $\id{U}$ or $\tape{c}$. We proceed with the case for $\tape{c}$, the one for $\id{U}$ is identical. Since $\tape{c}\colon U\to V_j$ then $\t_4$ should have type $V_j \to \bigoplus_{i=1}^n V_i$. This entails that $\t_4= \cobang{\bigoplus_{i=1}^{j-1} V_i} \oplus \id{V_j} \oplus \cobang{\bigoplus_{i=j+1}^{n} V_i}$. Thus $\t= \tape{c};\cobang{\bigoplus_{i=1}^{j-1} V_i} \oplus \id{V_j} \oplus \cobang{\bigoplus_{i=j+1}^{n} V_i}$ which, by the laws of symmetric monoidal categories, is equal to $\cobang{\bigoplus_{i=1}^{j-1} V_i} \oplus \tape{c} \oplus \cobang{\bigoplus_{i=j+1}^{n} V_i}$. Thus $\t$ is in the second shape.
\end{itemize}

Let us consider now the case $\t_1 \neq \id{U}$. For any $j\in \{1,\dots, n\}$, $\t_4$ can be written as $\codiag{P}^x \oplus \codiag{Q}^{y}$ where 
$x,y\in \mathbb{N}$, $P=\bigoplus_{i=1}^{j} V_i$ and $Q=\bigoplus_{i=j+1}^{n} V_i$
and  $\codiag{P}^x\colon \bigoplus_{k=1}^x P \to P$ and $\codiag{Q}^y \colon  \bigoplus_{k=1}^yQ\to Q$ are the arrows defined in (\ref{eq:codiagn in appendice}).

Now $\t_3$ should have as codomain $ \bigoplus_{k=1}^x P \oplus \bigoplus_{k=1}^yQ$. Since $\t$ has domain $U$, $\t_1; \t_2$ has type $U \to \bigoplus_{i=1}^z U$ for some $z\in \mathbb{N}$. Since the arrows in $\t_3$ preserve the size of domain and codomains we have that, the domain of $\t_3$ can be written as $ \bigoplus_{k=1}^x(\bigoplus_{i=1}^{j} U) \oplus \bigoplus_{k=1}^y(\bigoplus_{i=j+1}^{n} U)$. This fact entails that $\t_3= \t_3^1 \oplus \t_3^2$ for some $\t_3^1 \colon  \bigoplus_{k=1}^x(\bigoplus_{i=1}^{j} U) \to  \bigoplus_{k=1}^x(\bigoplus_{i=1}^{j} V_i)$ and  $\t_3^2 \colon  \bigoplus_{k=1}^y(\bigoplus_{i=j+1}^{n} U) \to  \bigoplus_{k=1}^y(\bigoplus_{i=j+1}^{n} V_i)$. 

Now $\t_2$ should have as codomain $ \bigoplus_{k=1}^x(\bigoplus_{i=1}^{j} U) \oplus \bigoplus_{k=1}^y(\bigoplus_{i=j+1}^{n} U)$. Since the arrows in the shape of $\t_2$ can only decrease the size of codomain, $\t_2$ has domain $ \bigoplus_{k=1}^x(\bigoplus_{i=1}^{x'} U) \oplus \bigoplus_{k=1}^y(\bigoplus_{i=1}^{y'} U)$ for some $x'\geq j$ and $y'\geq n-j+1$.
This fact entails that $\t_2=\t_2^1 \oplus \t_2^2$ for some $\t_2^1 \colon \bigoplus_{k=1}^x(\bigoplus_{i=1}^{x'} U) \to  \bigoplus_{k=1}^x(\bigoplus_{i=1}^{j} U)$ and $\t_2^2 \colon \bigoplus_{k=1}^y(\bigoplus_{i=1}^{y'} U) \to  \bigoplus_{k=1}^y(\bigoplus_{i=j+1}^{n} U)$.

Now $\t_1$ has type $U \to \bigoplus_{k=1}^x(\bigoplus_{i=1}^{x'} U) \oplus \bigoplus_{k=1}^y(\bigoplus_{i=1}^{y'} U)$. Note that the latter object is different from $U$ as $\t_1\neq \id{U}$. By using \eqref{eq:diagp assoc} and \eqref{eq:diagp symmetry}, $\t_1$ can be written in the form $\diagp{}; (\t_1^1\oplus \t_1^2)$ for some $\t_1^1\colon U \to \bigoplus_{k=1}^x(\bigoplus_{i=1}^{x'} U)$ and $\t_1^2\colon U \to \bigoplus_{k=1}^y(\bigoplus_{i=1}^{y'} U)$.

In summary, $\t=\diagp{};(\, (\t_1^1;\t_2^1; \t_3^1; \codiag{P}^x) \oplus (\t_1^2;\t_2^2; \t_3^2; \codiag{Q}^y) \,)$ is in the third form.

\end{proof}
\begin{proof}[Proof of Lemma~\ref{decomposition}]
Let $\t\colon U \to \bigoplus_{i=1}^n Q_i$ be an arrow of $\CatTapeC$ and observe that every $Q_i=\bigoplus_{j=1}^{m_i} U_j^i$ for some $m_i\in \mathbb{N}$ and $U_j^i$ objects of $\Cat{C}$. Hence, $\bigoplus_{i=1}^n Q_i = \bigoplus_{i=1}^n \bigoplus_{j=1}^{m_i} U_j^i= \bigoplus_{l=1}^s V_l$ for $s=m_1+\dots + m_n$ and $V_l = U_j^i$ if $l=m_1+\dots + m_{i-1}+j$ and  $j\in \{1,\dots,m_i\}$. Now we proceed by induction on $n$. The case $n=0$ is trivial since $\bigoplus_{i=1}^n Q_i$ is the final object $\zero$ and $\t= \bangp{U}=\,\diagpn{\vec{p}\;}{U}{0}$. For the induction case $n+1$, consider $\t\colon U \to Q_1 \oplus \bigoplus_{j=1}^n R_j$ for $R_j=Q_{j+1}$. By Lemma~\ref{lemma:division}, $\t$ is in one of the forms indicated in the statement of Lemma~\ref{lemma:division}. If $\t= \bang{U}; \cobang{\bigoplus_{l=1}^s V_l}= \bang{U}; \cobang{\bigoplus_{i=1}^{n+1} Q_i}$, then taking all $p_i=0$ it follows that $\t=\diagpn{\;\vec{p}}{U}{n+1} ; \bigoplus_{i=1}^{n+1} \star_{U,Q_i}$. Now consider the case in which there exists $k\in \{1,\dots,s\}$ and $\t_k\colon U \to V_k$ such that $\t = \cobang{\bigoplus_{l=1}^{k-1} V_l} \oplus \t_k \oplus \cobang{\bigoplus_{l=k+1}^{s} V_l}$. Assume that $k=m_1+\dots + m_{h-1}+j$ for $j\in \{1,\dots,m_h\}$ (then $V_k= U_j^h$). Hence, $\t= \cobang{\bigoplus_{i=1}^{h-1} Q_i} \oplus \t' \oplus \cobang{\bigoplus_{i=h+1}^{n+1} Q_i}$ where $\t' = \cobang{\bigoplus_{i=1}^{j-1} U_i^h} \oplus \t_k \oplus \cobang{\bigoplus_{i=j+1}^{m_h} U_i^h}$. Taking $p_i=0$ for $i\neq h$ and $p_h=1$, it follows that $\t=\diagpn{\;\vec{p}}{U}{n+1} ; \bigoplus_{i=1}^{h-1} \star_{U,Q_i} \oplus \t' \oplus \bigoplus_{i=h+1}^{n+1} \star_{U,Q_i}$. Now consider the last case and assume that $\t=\,\diagp{U};(\s_1\oplus \s_2)$, with $\s_1\colon U \to Q_1$ and $\s_2\colon U \to \bigoplus_{j=1}^n R_j$. By induction hypothesis, $\s_2$ can be written as $\s_2 = \diagpn{\;\vec{p}}{U}{n} ; \bigoplus_{i=1}^n \s_2^j$ where $\s_2^j\colon U \to Q_{j+1}$ for $j=1,\dots, n$. Hence, $\t =\, \diagp{U};(\s_1\oplus \s_2) = \diagp{U};(\id{U} \oplus \diagpn{\;\vec{p}}{U}{n}) ; (\s_1 \oplus \bigoplus_{i=1}^n \s_2^j)$. Taking $t_1=p$ and $t_i = (1-p)\cdot p_{i-1}$ for $i=2,\dots, n+1$, it follows that $\t=\diagpn{\;\vec{t}}{U}{n+1} ; \bigoplus_{i=1}^{n+1}\s'_i$, where $\s'_1 = \s_1$ and $\s'_i = \s_2^{i-1}$ for $i=2,\dots, n+1$.

%If $\t= \bang{U}; \cobang{\bigoplus_{i=1}^n U_i}$, then take all $p_i=0$ and $\t_i = \cobang{U_i}$. Simple computations proves that $\t=\diagpn{\;\vec{p}}{U}{n} ; \bigoplus_{i=1}^n \t_i$. 
%
%If $\t = \cobang{\bigoplus_{i=1}^{j-1} U_i} \oplus \t_j \oplus \cobang{\bigoplus_{i=j+1}^{n} U_i}$, then take $p_j=1$ and, for $i\neq j$, $p_i = 0$ and $\t_i =\cobang{U_i}$. Simple computations proves that $\t=\diagpn{\;\vec{p}}{U}{n} ; \bigoplus_{i=1}^n \t_i$.
%
%In the third case, take $j=n$.
\end{proof}



\subsection{Cancellativity}\label{app:cancellativity}
In order to prove cancellativity, we first consider the case of arrows of type $U \to V$ where $U,V$ are objects of $\Cat{C}$.
\begin{lemma}\label{lemma:cancellativity11}
For all $r\in (0,1)$, for all $\s,\t \colon U \to V$, 
if $r\cdot \s = r\cdot \t$ then $\s = \t$.
\end{lemma}
\begin{proof}
By Lemma \ref{lemma: arrows A-B subdistributions}, we know $\CatTapeC[U,V]= \subdistr(\Cat{C}[U,V])$. It is well known (see e.g. \cite{sokolova2018termination}) that, for all sets $X$, $\subdistr(X)$ is a cancellative pca.
\end{proof}

Then, we consider the case of arrows of type $U \to \bigoplus_{i=1}^n V_i$ where $U,V_i$ are objects of $\Cat{C}$.
\begin{lemma}\label{lemma:cancellativityhalf}
For all $n \in \mathbb{N}$, for all $r\in (0,1)$, for all $\s,\t \colon U \to \bigoplus_{i=1}^n V_i$, 
if $r\cdot \s = r\cdot \t$ then $\s = \t$.
\end{lemma}
\begin{proof}
We proceed by induction on $n$. 

The base case $n=0$ is trivial since $\bigoplus_{i=1}^n V_i$ is the final object $\zero$.

For the induction case $n+1$, we exploit Lemma~\ref{lemma:division} and we consider the case where
\begin{equation}\label{eq:st}
\s=\, \diagq{} ; (\s_1 \oplus \s_2)   \qquad \text{ and } \qquad \t =\, \diagp{}; (\t_1 \oplus \t_2)
\end{equation}
for $\s_1, \t_1 \colon U \to \bigoplus_{i=1}^n V_i$ and $\s_2,\t_2 \colon U \to V_{n+1}$. The other cases are trivial.

By Lemma \ref{lemma:basicTC}.3, we have that
\begin{equation}\label{eq:simple}
r\cdot \s=\, \diagq{} ; (r\cdot \s_1 \oplus r\cdot \s_2)   \qquad \text{ and } \qquad r\cdot  \t =\, \diagp{}; (r\cdot \t_1 \oplus r\cdot \t_2)
\end{equation}
Since by hypothesis $r\cdot \s = r \cdot \t$, it holds both  
\[r\cdot \s ; (\id{} \oplus \bang{}) = r \cdot \t ; (\id{} \oplus \bang{}) \qquad \text{ and } \qquad r\cdot \s ; ( \bang{} \oplus \id{}) = r \cdot \t ;( \bang{} \oplus \id{})\text{.}\]
Thus, by \eqref{eq:simple} and Lemma~\ref{lemma:basicTC}.1 and 2, one obtains that
\begin{equation}\label{eq:twoeq}q\cdot r \cdot \s_1 = p \cdot r \cdot \t_1 \qquad \text{ and } \qquad (1-q) \cdot r \cdot \s_2 = (1-p) \cdot r \cdot \t_2\text{.}\end{equation}

There are now three cases: $p=q$, $p>q$ and $p<q$.

If $p=q$, then  $\s_1=\t_1$ by induction hypothesis and  $\s_2=\t_2$ by Lemma~\ref{lemma:cancellativity11}. Then, by \eqref{eq:st}, $\s=\t$.

If $p>q$, we can rewrite \eqref{eq:twoeq} as
\begin{equation*}p\cdot r \cdot \frac{q}{p} \cdot \s_1 = p \cdot r \cdot \t_1 \qquad \text{ and } \qquad (1-q) \cdot r \cdot \s_2 = (1-q) \cdot r \cdot \frac{1-p}{1-q} \cdot \t_2\text{.}\end{equation*}
Thus, by induction hypothesis $ \frac{q}{p} \cdot \s_1 = \t_1$ and, by Lemma~\ref{lemma:cancellativity11}, $\s_2 =\frac{1-p}{1-q} \cdot \t_2$. By \eqref{eq:st} and Lemma~\ref{lemma:fractions}, we have that $\s=\t$.
The case $p<q$ is symmetrical.
\end{proof}

We can finally consider the general case.

\begin{proof}[Proof of Lemma~\ref{cancellativity}]
Let $P=\bigoplus_{i=1}^n U_i$. We proceed by induction on $n$. The base case $n=0$ is trivial since $\bigoplus_{i=1}^n U_i$ is the initial object $\zero$.

For the induction case $n+1$, since $\CatTapeC$ has finite coproducts, one can assume without loss of generality that
\begin{equation}\label{eq:s1s2coproduct}
\s = (\s_1 \oplus \s_2) ; \codiag{Q} \qquad \text{ and } \t = (\t_1 \oplus \t_2) ; \codiag{Q}
\end{equation}
for some $\s_2,\t_2 \colon U_{n+1} \to Q$ and  $\s_1,\t_1\colon P' \to Q$ where $P'=\bigoplus_{i=1}^n U_i$. Note moreover that the coproduct universal property entails that
\begin{equation}\label{eq:ifflocal}
\s=\t \text{ iff } \s_1=\t_1 \text{ and } \s_2=\t_2\text{.}
\end{equation}

By Lemma~\ref{lemma:pcopairing}, we have that
$r \cdot \s = (r\cdot \s_1 ) \oplus (r\cdot  \s_2 ); \codiag{Q}$ and $r \cdot \t = (r\cdot \t_1 ) \oplus (r\cdot  \t_2 ); \codiag{Q} $. Thus, by \eqref{eq:ifflocal} we have that
 \[r\cdot \s_1= r\cdot \t_1 \qquad \text{ and }\qquad r\cdot \s_2= r\cdot \t_2 \]
 From the rightmost equality above and Lemma~\ref{lemma:cancellativityhalf}, we derive that $\s_2=\t_2$. From the leftmost equality above and induction hypothesis, we derive that $\s_1=\t_1$. By \eqref{eq:ifflocal}, it holds that $\s=\t$.
 \end{proof}



%%%%%%CAMBIO NOME A QUESTA APPENDICE DA CONVEX PRODUCTS A JOINTLY MONIC PROJECTIONS%%%%%%%%%
\subsection{Jointly Monic Projections}
%Now, proving $<\t,\s>_{p,1-p}$ is the unique mediating morphism is pretty easy in the case $p\in \{0,1\}$. It follows immediately by the following results.



\begin{lemma}\label{keyLemma001}
Let $\t\colon U \to Q \oplus R$ be an arrow in $\CatTapeC$. 
\begin{enumerate} 
\item If $\t; (\id{Q} \oplus\, \bangp{R}) =\s$
and $\t; (\,\bang{Q} \oplus \id{R}) = \bang{U}; \cobang{R}$ then $\t =\s \oplus \cobang{R}$;
\item If $\t; (\id{Q} \oplus\, \bangp{R}) =\bang{U}; \cobang{Q}$  
and $\t; (\,\bang{Q} \oplus \id{R}) = \s$ then $\t = \cobang{Q} \oplus \s$;
\end{enumerate}
\end{lemma}
\begin{proof}
We prove the first point: the second is symmetrical.

First, observe that $\t$ should be in one of the forms of Lemma~\ref{lemma:division}. The only challenging case is the third form (iii) $\t=\,\diagp{U}; (\t_1 \oplus \t_2)$ for some $\t_1\colon U \to Q$ and $\t_2\colon U \to R$. Since
\begin{align*}
\bang{U}; \cobang{R} &= \t; (\,\bang{Q} \oplus \id{R}) \tag{Hypothesis}\\
&=\, \diagp{U}; (\t_1 \oplus \t_2) ; (\,\bang{Q} \oplus \id{R}) \tag{iii}\\
&=\, \diagp{U}; ((\t_1 ; \bang{Q}) \oplus \t_2) \tag{Symmetric Monoidal Category}\\
&=\, \diagp{U} ; (\bang{U} \oplus \t_2) \tag{\ref{eq:bangp nat}}\\
&= (1-p)\cdot \t_2 \tag{Lemma \ref{lemma:basicTC}.2}
\end{align*}
%\marginpar{qui al posto di dire by lemma 41, ho usato cancellatività}then %by Lemma \ref{lemma: arrows A-B subdistributions}, $\t_2 = \bang{U}; \cobang{R}$.
since also $(1-p)\cdot (\bang{U}; \cobang{R})=\bang{U}; \cobang{R}$, by cancellativity (Lemma~\ref{cancellativity}), we have that $\t_2 = \bang{U}; \cobang{R}$.
Thus
\begin{align*}
\s \oplus \cobang{R} & = (\t; (\id{Q} \oplus\, \bangp{R})) \oplus \cobang{R} \tag{Hypothesis}\\
&= (\diagp{U}; (\t_1 \oplus \t_2) ; (\id{Q} \oplus\, \bangp{R})) \oplus \cobang{R} \tag{iii}\\
&= (\diagp{U}; (\t_1 \oplus (\t_2 ; \bang{R})))  \oplus\, \cobang{R} \tag{Symmetric Monoidal Category}\\
&= (\diagp{U}; (\t_1 \oplus ( \bang{U}; \cobang{R} ; \bang{R})))  \oplus\, \cobang{R} \tag{$\t_2 = \bang{U}; \cobang{R}$}\\
&= (\diagp{U}; (\t_1 \oplus  \bang{U}) ) \oplus \cobang{R} \tag{\ref{eq:bangp nat}}\\
&=\, \diagp{U}; (\t_1 \oplus  (\bang{U} ; \cobang{R}))\tag{Symmetric Monoidal Category}\\
&=\, \diagp{U}; (\t_1 \oplus  \t_2)\tag{$\t_2 = \bang{U}; \cobang{R}$}\\
&= \t \tag{iii}\\
\end{align*}
\end{proof}


%%%%%%%%%%%%%%LA VERSIONE GENRALE DEL LEMMA PRECEDENTE NON CI SERVE PIù%%%%%%%%%%%%%%%%%%%%
\begin{comment}

The lemma above is generalised to arbitrary arrows as follows. The proof substantially relies on the fact that $\CatTapeC$ has coproducts.


\begin{lemma}\label{keyLemma01}
Let  $\t\colon P \to Q \oplus R$ be an arrow in $\CatTapeC$ where $P=\bigoplus_{i=1}^nU_i$. 
\begin{enumerate} 
\item If $\t; (\id{Q} \oplus\, \bangp{R}) =\s$
and $\t; (\,\bang{Q} \oplus \id{R}) = \bang{P}; \cobang{R}$ then $\t =\s \oplus \cobang{R}$;
\item If $\t; (\id{Q} \oplus\, \bangp{R}) =\bang{P}; \cobang{Q}$  
and $\t; (\,\bang{Q} \oplus \id{R}) = \s$ then $\t = \cobang{Q} \oplus \s$;
\end{enumerate}
\end{lemma}
\begin{proof}
    The proof proceeds by induction on $n$. The base case $n=0$ is trivial since $\bigoplus_{i=1}^0 U_i$ is the initial object $\zero$ .
  %  For the induction case $n+1$, since $\CatTapeC$ has finite coproducts, then \begin{equation} \label{eq:local11}
%\t_1 = (\s_1^1 \oplus \s_1^2) ; \codiag{Q} \qquad \text{ and } \qquad \t_2 = (\s_2^1 \oplus \s_2^2) ; \codiag{R}
%\end{equation}
%for $\s_1^1 \colon  \bigoplus_{i=1}^nU_i \to Q$, $\s_1^2\colon U_{n+1} \to Q$, $\s_2^1 \colon  \bigoplus_{i=1}^nU_i \to R$, $\s_2^2\colon U_{n+1} \to R$.

For the induction case $n+1$, since $\CatTapeC$ has finite coproducts, we can assume that $\t= (\t_1 \oplus \t_2) ; \codiag{Q\piu R}$ where $\t_1\colon \bigoplus_{i=1}^nU_i\to Q\piu R$ 
and $\t_2\colon U_{n+1} \to Q\piu R$ and $\s= (\s_1 \oplus \s_2) ; \codiag{Q}$ where $\s_1\colon \bigoplus_{i=1}^nU_i\to Q$ and $\s_2\colon U_{n+1} \to Q$.

Now, if $\t; (\id{Q} \oplus\, \bangp{R}) =\s$ it follows that 
\begin{align} \t_1\piu \t_2 ; \codiag{Q\piu R} ; (\id{Q} \oplus\, \bangp{R}) &= \t_1\piu \t_2 ; (\id{Q} \oplus\, \bangp{R}) \piu (\id{Q} \oplus\, \bangp{R}) ; \codiag{Q} \tag{axiom \ref{eq:codiag nat}}
\end{align}
and hence $\t_1;(\id{Q} \oplus\, \bangp{R}) = \s_1$ and $\t_2;(\id{Q} \oplus\, \bangp{R}) = \s_2$. Similarly, from $\t; (\,\bang{Q} \oplus \id{R}) = \bang{P}; \cobang{R}$, using the axiom \ref{eq:codiag nat} and \ref{eq:bangp coherence}, we have that $\t_1;(\bang{Q} \oplus \id{R}) = \bang{\bigoplus_{i=1}^nU_i}; \cobang{Q}$ and $\t_2;(\bang{Q} \oplus \id{R}) = \bang{U_{n+1}}; \cobang{R}$. Hence, by inductive hypothesis, we have that $\t_1 = \s_1 \oplus \cobang{R}$ and, from Lemma \ref{keyLemma001}, $\t_2 = \s_2 \oplus \cobang{R}$. Now, the statement is obtained as follows
\begin{align}
    \t_1\piu \t_2 ; \codiag{Q\piu R} &= (\s_1 \oplus \cobang{R}) \piu (\s_2 \oplus \cobang{R}) ; \codiag{Q\piu R} \tag*{}\\
    &= (\s_1 \oplus \cobang{R}) \piu (\s_2 \oplus \cobang{R})  ; (\id{Q} \piu \sigma_{R,Q} \piu \id{R}) ; \codiag{Q}\piu \codiag{R} \tag{axiom \ref{eq:codiag coherence}}\\
    &= (\s_1\oplus \s_2) ; \codiag{Q}\piu \cobang{R} \tag{axiom \ref{eq:codiag nat}}\\
    &=\s\piu \cobang{R} \notag
\end{align}

\end{proof}

\end{comment}

%Proving that $<\t_1,\t_2>_{p,1-p}$ is the convex pairing in the case of $p\in(0,1)$ relies on cancellativity. 
%We first prove the case for $\t_1\colon U\to P$ and $\t_2\colon U \to Q$ where $U$ is an object of $\Cat{C}$

\begin{lemma}\label{lemma:comvexproductpreliminary}
Let $\t_1\colon U\to P$ and $\t_2\colon U \to Q$ be arrows of $\CatTapeC$ and $p\in (0,1)$. Then $\diagp{U}; (\t_1 \oplus \t_2) \colon U \to P\oplus Q$ is the unique arrow $\t$ such that
$p\cdot \t_1 = \t; (\id{P} \oplus \bang{Q})$ and $(1-p)\cdot \t_2 = \t; (\bang{P} \oplus \id{Q})$.
\end{lemma}
\begin{proof}
First observe that by Lemma~\ref{lemma:basicTC}.1 and 2 we have that $\diagp{U}; (\t_1 \oplus \t_2) ; (\id{P} \oplus \bang{Q}) = p\cdot \t_1$ and that $\diagp{U}; (\t_1 \oplus \t_2) ; (\bang{P} \oplus \id{Q}) = (1-p) \cdot \t_2$.

We now need to prove that $\diagp{U}; (\t_1 \oplus \t_2)$ is the unique arrow with such property.  We prove that if (a) $p\cdot \t_1 = \t; (\id{P} \oplus \bang{Q})$ and (b) $(1-p)\cdot \t_2 = \t; (\bang{P} \oplus \id{Q})$ then $\t =\, \diagp{U}; (\t_1 \oplus \t_2)$.

Note that $\t$ can be in one of the forms of Lemma~\ref{lemma:division}. The only relevant case is the last one: $\t =\, \diagpX{p'}{U} ; (\t_1' \oplus 	\t_2')$. By (a), (b) Lemma~\ref{lemma:basicTC}.1 and 2, it holds that 
\begin{center}$p\cdot \t_1 = p' \cdot \t_1'$ and  $(1-p)\cdot \t_2 = (1-p') \cdot \t_2'$. \end{center}

There are now three cases:
\begin{itemize}
\item If $p=p'$, then by Lemma~\ref{cancellativity}, $\t_1=\t_1'$ and $\t_2=\t_2'$. Thus, $\t =\, \diagp{U}; (\t_1 \oplus \t_2) $.
\item If $p>p'$, then we can rewrite the above equalities as  $p\cdot \t_1 = p \cdot \frac{p'}{p} \cdot \t_1'$ and $(1-p')\cdot\frac{(1-p)}{1-p'}\cdot \t_2 = (1-p') \cdot \t_2'$. By Lemma~\ref{cancellativity}, we have that $\t_1=\frac{p'}{p}\cdot\t_1'$ and $\frac{1-p}{1-p'} \cdot \t_2 =\t_2'$. Thus, by Lemma~\ref{lemma:fractions}, $\t=\, \diagp{U}; (\t_1 \oplus \t_2)$.
\item The case $p<p'$ is symmetrical to the one above.
\end{itemize}
\end{proof}

\begin{corollary}\label{cor:finaleTC}
    Let $\t_1\colon U\to P$ and $\t_2\colon U \to Q$ be arrows of $\CatTapeC$ and $p,q\in [0,1]$ such that $p+q\le 1$. Then $\diaggen{p,q}{U};(\t_1 \oplus \t_2) \colon U \to P\oplus Q$ is the unique arrow $\t$ such that  $p\cdot \t_1 = \t; (\id{P} \oplus \bang{Q})$ and $q\cdot \t_2 = \t; (\bang{P} \oplus \id{Q})$.
\end{corollary}
\begin{proof}
    The case where $p$ or $q$ is equal to $1$ follows from Lemma~\ref{keyLemma001} and the fact that $\diaggen{1,0}{U} =\, \diaggen{1}{U}= \id{U}\oplus \cobang{U}$ and $\diaggen{0,1}{U} =\, \diaggen{0}{U} = \cobang{U}\oplus \id{U}$. The case $p,q\in (0,1)$ and $p+q=1$ follows from Lemma~\ref{lemma:comvexproductpreliminary} and the fact that $\diaggen{p,q}{U} =\, \diagp{U}$. The case where $p,q\in (0,1)$ and $p+q<1$ follows from the previous case and the fact that $\diaggen{p,q}{U};(\t_1\oplus \t_2) =\, \diaggen{p,1-p}{U}; (\t_1 \oplus \frac{q}{1-p}\cdot \t_2)$.
\end{proof}


%%%%%%%%COMMENTO LA VERSIONE GENERALE PERCHé NON LA USIAMO PIU%%%%%

\begin{comment}
    

The lemma above is generalised to arbitrary arrows as follows. The proof substantially relies on the fact that $\CatTapeC$ has coproducts.

\begin{lemma}\label{lemma:comvexproductalmost}
Let $\t_1\colon \bigoplus_{i=1}^nU_i\to P$ and $\t_2\colon \bigoplus_{i=1}^nU_i \to Q$ be arrows of $\CatTapeC$ and $p\in (0,1)$. Then there exists a unique arrow $\t \colon \bigoplus_{i=1}^n U_i \to P\oplus Q$ such that
$p\cdot \t_1 = \t; (\id{P} \oplus \bang{Q})$ and $(1-p)\cdot \t_2 = \t; (\bang{P} \oplus \id{Q})$.
\end{lemma}
\begin{proof}
The proof proceed by induction on $n$. The case $n=0$ is trivial since $\bigoplus_{i=1}^0U_i$ is the initial object $\zero$.

We consider the case $n+1$. Since $\CatTapeC$ has coproducts then 
\begin{equation} \label{eq:local2}
\t_1 = (\s_1^1 \oplus \s_1^2) ; \codiag{P} \qquad \text{ and } \qquad \t_2 = (\s_2^1 \oplus \s_2^2) ; \codiag{Q}
\end{equation}
for $\s_1^1 \colon  \bigoplus_{i=1}^nU_i \to P$, $\s_1^2\colon U_{n+1} \to P$, $\s_2^1 \colon  \bigoplus_{i=1}^nU_i \to Q$, $\s_2^2\colon U_{n+1} \to Q$.

By induction hypothesis there exists a unique arrow $\t \colon  \bigoplus_{i=1}^nU_i \to P \oplus Q$ such that
\begin{equation}\label{eq:local3}
\t; (\id{P} \oplus \bang{Q}) = p\cdot \s_1^1 \qquad \text{ and } \qquad \t; (\bang{P} \oplus \id{Q}) = (1-p)\cdot \s_2^1
\end{equation}

We now focus on  $\t \oplus ( \diagp{U_{n+1}} ; (\s_1^2 \oplus \s_2^2) ); \codiag{P\oplus Q} \colon \bigoplus_{i=1}^{n+1} U_i \to P\oplus Q$.
\begin{align*}
 & \t \oplus ( \diagp{U_{n+1}} ; (\s_1^2 \oplus \s_2^2) ); \codiag{P\oplus Q} ; (\id{P} \oplus \bang{Q}) \\& = ( \, (\t; (\id{P} \oplus \bang{Q})) \oplus (\diagp{U_{n+1}} ; (\s_1^2 \oplus \s_2^2) ; (\id{P} \oplus \bang{Q})) \,) ; \codiag{P\oplus Q} \tag{nat. of $\codiag{P\oplus Q}$}\\
&= ( \, (\t; (\id{P} \oplus \bang{Q})) \oplus (\diagp{U_{n+1}} ; (\s_1^2 \oplus \bang{U_{n+1}})) \,) ; \codiag{P\oplus Q} \tag{nat. of $\bang{}$}\\
&= ( \, (\t; (\id{P} \oplus \bang{Q})) \oplus p\cdot \s_1^2  \,) ; \codiag{P\oplus Q} \tag{Lemma \ref{lemma:pcdot}}\\
&= ( \, p\cdot \s_1^1 \oplus p\cdot \s_1^2  \,) ; \codiag{P\oplus Q} \tag{\ref{eq:local3}}\\
&=  p\cdot  ( \, ( \s_1^1 \oplus \s_1^2 ) ; \codiag{P\oplus Q} \, ) \tag{Lemma \ref{lemma:pcopairing}}\\
&=  p\cdot  \t_1 \tag{\ref{eq:local2}}\\
\end{align*}
Similar computations proves that $\t \oplus ( \diagp{U_{n+1}} ; (\s_1^2 \oplus \s_2^2) ); \codiag{P\oplus Q} ; (\bang{P} \oplus \id{Q}) = (1-p) \cdot \t_2$.

Now we need to prove that $\t \oplus ( \diagp{U_{n+1}} ; (\s_1^2 \oplus \s_2^2) ); \codiag{P\oplus Q}$ is the unique arrows with such property.
Assume that $\t'\colon \bigoplus_{i=1}^{n+1} U_i \to P\oplus Q$ is such that
\begin{equation}\label{eq:local4}
\t'; (\id{P} \oplus \bang{Q}) = p\cdot  \t_1 \qquad \text{ and }\qquad \t'; (\bang{P} \oplus \id{Q}) = (1-p)\cdot  \t_2
\end{equation}
Since $\CatTapeC$ has coproducts $\t' = (\t_1' \oplus \t_2' ); \codiag{P\oplus Q}$ for some $\t_1'\colon \bigoplus_{i=1}^{n} U_i \to P\oplus Q $ and $\t_2'\colon U_{n+1} \to P \oplus Q$.
 By using naturality of $\codiag{P\oplus Q}$ one readily has
 \begin{equation}\label{eq:local5}
 \t' ; (\id{P} \oplus \bang{Q}) = (\, (\t_1'; (\id{P} \oplus \bang{Q})) \oplus (\t_2' ; (\id{P} \oplus \bang{Q})) \,); \codiag{P\oplus Q}\text{.}
 \end{equation}
% and
% \begin{equation*}
% \t' ; (\bang{P} \oplus \id{Q}) = (\, (\t_1'; (\bang{P} \oplus \id{Q})) \oplus (\t_2' ; (\bang{P} \oplus \id{Q})) \,); \codiag{P\oplus Q}
% \end{equation*}
 Now, it holds that
 \begin{align*}
  (\, (\t_1'; (\id{P} \oplus \bang{Q})) \oplus (\t_2' ; (\id{P} \oplus \bang{Q})) \,); \codiag{P\oplus Q} & = \t' ; (\id{P} \oplus \bang{Q}) \tag{\ref{eq:local5}} \\
 &= p\cdot  \t_1 \tag{\ref{eq:local4}}\\
 &= ( \, p\cdot \s_1^1 \oplus p\cdot \s_1^2  \,) ; \codiag{P\oplus Q} \tag{derivation above}
 \end{align*}
 and thus, by the universal property of coproducts, we can conclude that
 \begin{equation}\label{eq:local6}
  \t_1'; (\id{P} \oplus \bang{Q}) = p\cdot \s_1^1 \qquad \text{ and } \qquad \t_2' ; (\id{P} \oplus \bang{Q}) = p\cdot \s_1^2\text{.}
 \end{equation}
 With the symmetric argument, one obtains that
 \begin{equation}\label{eq:local7}
  \t_1'; (\bang{P} \oplus \id{Q}) = (1-p)\cdot \s_2^1 \qquad \text{ and } \qquad \t_2' ; (\bang{P} \oplus \id{Q}) = (1-p)\cdot \s_2^2\text{.}
 \end{equation}
Now, by the leftmost equations in \eqref{eq:local6} and \eqref{eq:local7} and the fact that $\t$ is the unique satisfying \eqref{eq:local3}, we obtain that $\t_1' = t$.

By Lemma~\ref{lemma:comvexproductpreliminary}, $\diagp{U_{n+1}} ; (\s_1^2 \oplus \s_2^2)$ is the unique arrow $h\colon  U_{n+1} \to P \oplus Q $ such that $h; (\id{P} \oplus \bang{Q}) = p\cdot \s_1^2$ and  $h; (\bang{P} \oplus \id{Q}) = (1-p)\cdot \s_2^2$. Thus, the rightmost equations in \eqref{eq:local6} and \eqref{eq:local7} entails that $\t_2'= \diagp{U_{n+1}} ; (\s_1^2 \oplus \s_2^2)$. This proves that $\t'= \t \oplus ( \diagp{U_{n+1}} ; (\s_1^2 \oplus \s_2^2) ); \codiag{P\oplus Q}$.
\end{proof}
\end{comment}


%%%%%%%%%%%COMMENTO PROP TC CONVEX PRODUCT%%%%%%%%%%%%%%%%%%
\begin{comment}
    

We can finally state our main result.
\begin{proposition}\label{prop:TChasconvexcoprduct}
$\CatTapeC$ has convex products.
\end{proposition}
\begin{proof}
We need to prove that for all arrows $\t_1\colon R\to P$ and $\t_2\colon R \to Q$ and for all $p,q\in[0,1]$ such that $p+q\leq 1$, there exists a unique $\t\colon R \to P\oplus Q$ such that $\t;\pi_1 =p\cdot \t_1$ and $\t;\pi_2 = q \cdot \t_2$.

The case when either $p=1$ or $q=1$ is proved by Lemma~\ref{keyLemma01}.
The case when $p,q\in (0,1)$ and $p+q=1$ is proved by Lemma~\ref{lemma:comvexproductalmost}.
The case when $p,q\in (0,1)$ and $p+q<1$ easily follows from the previous case: if $p\geq 1-q$ one can take $\langle f,g \rangle_{p,q} = \langle f, \frac{q}{1-p} \cdot g \rangle_{p,1-p}$.
If $p\leq 1-q$ one can take $\langle f,g \rangle_{p,q} = \langle \frac{p}{1-q} \cdot f, g \rangle_{1-q,q}$.
\end{proof}
\end{comment}
%%%%%%%%%%%END COMMENTO PROP TC CONVEX PRODUCT%%%%%%%%%%%%%%%%%%



\begin{proof}[Proof of Lemma~\ref{lemma:TC jmono}]
     First, we consider two arrows $\t,\s\colon U \to Q\oplus R$ where $U$ is an object of $\Cat{C}$ and assume that $\t;\pi_1 = \s;\pi_1$ and $\t;\pi_2 = \s;\pi_2$. By Lemma~\ref{decomposition}, there exist $\t_1,\s_1\colon U \to Q$, $\t_2,\s_2\colon U \to R$ and $p,q,p',q'\in [0,1]$ where $p+q\le 1$ and $p'+q'\le 1$ such that
     \[\t=\, \diaggen{p,q}{};(\t_1 \oplus \t_2) \text{ and }\s=\,\diaggen{p',q'}{};(\s_1 \oplus \s_2)\text{.}\]
The last two items of Lemma~\ref{lemma:diagpq} guarantee that $\t;\pi_1 = p\cdot \t_1$ and $\t;\pi_2 = q\cdot \t_2$ and $\s;\pi_1 = p'\cdot \s_1$ and $\s;\pi_2 = q'\cdot \s_2$. Since, by hypothesis, $\t;\pi_1 = \s;\pi_1$ and $\t;\pi_2 = \s;\pi_2$, we have that $p\cdot \t_1 = p'\cdot \s_1$ and $q\cdot \t_2 = q'\cdot \s_2$. 
    Finally, Corollary~\ref{cor:finaleTC} implies that $\t=\s$.
    
    
    For the general case $\t,\s\colon P\to Q\oplus R$ where $P=\bigoplus_{i=1}^n U_i$, we rely on the fact that $\CatTapeC$ has coproducts. Indeed, by the universal property of coproducts, we can write $\t = [\t_1,\dots,\t_n]$ and $\s = [\s_1,\dots,\s_n]$ for some $\t_i,\s_i\colon U_i \to Q\oplus R$. Now, from $\t;\pi_1 = \s;\pi_1$ and $\t;\pi_2 = \s;\pi_2$ it follows that $\t_i;\pi_1 = \s_i;\pi_1$ and $\t_i;\pi_2 = \s_i;\pi_2$ for all $i=1,\dots,n$. By the previous case, we have that $\t_i=\s_i$ for all $i=1,\dots,n$ and hence $\t=\s$.
\end{proof}



\subsection{Proof of Proposition~\ref{cor: quotient category is convex biproduct }}
We conclude this appendix with a proof of Proposition~\ref{cor: quotient category is convex biproduct } that turns out to be useful when considering diagrammatic languages (see Section~\ref{sec:probbooltapes}). First, we need the following result.

\begin{lemma}\label{lemma: convex products quoziente}
Let $\Cat{C}$ be a category and let $\sim$  be a congruence relation on $\CatTapeC$ (w.r.t. $;$ and $\oplus$) that is cancellative: for every $p\in (0,1)$, if $p\cdot f\sim p\cdot g$ then $f\sim g$. For all $f_1,g_1\colon Z\to X_1$, $f_2,g_2\colon Z\to X_2$ and $p_1,p_2,q_1,q_2\in[0,1]$ such that $p_1+p_2\le 1$ and $q_1+q_2\le 1$
\begin{center}if $p_1\cdot f_1 \sim q_1\cdot g_1$ and $p_2\cdot f_2\sim q_2\cdot g_2$  then $\langle f_1,f_2\rangle_{(p_1,p_2)}\sim\langle g_1,g_2\rangle_{(q_1,q_2)}$.\end{center}
\end{lemma}

\begin{proof}\label{proof-lemma: convex products quoziente}%\marginpar{NEW-OK}
We first consider the case where $p_1,p_2,q_1,q_2\in(0,1)$. We assume that $p_1\leq q_1$  and $q_2\leq p_2$.  The cases of the other inequalities are proved similarly.

First observe that
\begin{align*}
(1-p_1)\cdot (\frac{p_2}{1-p_1}\cdot f_2) &= p_2\cdot f_2 \notag \\
& \sim q_2\cdot g_2 \tag{Hypothesis}\\
&= (1-p_1)\cdot (\frac{q_2}{1-p_1}\cdot g_2) \notag
\end{align*}
Thus, by cancellativity of $\sim$, it holds that $\frac{p_2}{1-p_1}\cdot f_2 \sim \frac{q_2}{1-p_1}\cdot g_2$ which, by Lemma \ref{lemma:initialproperties}, means
\begin{equation}\tag{$\dagger$}\label{eq:internal1}
\diagpX{\frac{p_2}{1-p_1}\,}{};(f_2\piu\, \bangp{}) \sim \diagpX{\frac{q_2}{1-p_1}\,}{};(g_2\piu\, \bangp{}) 
\end{equation}
One can similarly prove that 
\begin{equation}\tag{$\ddagger$}\label{eq:internal2}
\diagpX{\frac{p_1}{1-q_2}\,}{};(f_1\piu\, \bangp{}) \sim \diagpX{\frac{q_1}{1-q_2}\,}{};(g_1\piu\, \bangp{}) 
\end{equation}

We conclude with the following derivation. %The first equality is obtained as follows:
  \begin{align*}
        \langle f_1,f_2\rangle_{(p_1,p_2)}&=\, \diagpX{p_1}{}; (f_1\piu\, (\diagpX{\frac{p_2}{1-p_1}\,}{};(f_2\piu\, \bangp{})) ) \tag{Lemma~\ref{lemma:diagpq3}}\\
        &\sim \, \diagpX{p_1}{}; (f_1\piu\, (\diagpX{\frac{q_2}{1-p_1}\,}{};(g_2\piu\, \bangp{})) ) \tag{\eqref{eq:internal1}}\\
& = \, \diagpX{q_2}{};(g_2\piu\, (\diagpX{\frac{p_1}{1-q_2}\,}{};(f_1\piu\, \bangp{})) ) \tag{\ref{eq:diagp assoc} + \ref{eq:diagp symmetry}}  \\
&\sim \, \diagpX{q_2}{};(g_2\piu\, (\diagpX{\frac{q_1}{1-q_2}\,}{};(g_1\piu\, \bangp{})) ) \tag{\eqref{eq:internal2}}\\
 &=\,  \diagpX{q_1}{}; (g_1\piu\, (\diagpX{\frac{q_2}{1-q_1}\,}{};(g_2\piu\, \bangp{})) ) \tag{\ref{eq:diagp assoc} + \ref{eq:diagp symmetry}}\\
 &=  \langle g_1,g_2\rangle_{(q_1,q_2)}  \tag{Lemma~\ref{lemma:diagpq3}}
  \end{align*}
  
  
We now prove the case where $p_1=1$ and $p_2=0$. We have two sub-cases: either $q_1=1$ or $q_1\neq 1$. 
 If $q_1=1$, then by Lemma~\ref{lemma:diagpq3} easily follows that $\langle f_1,f_2\rangle_{(1,0)}=f_1;\iota_1$ and $\langle g_1,g_2\rangle_{(1,0)}=g_1;\iota_1$. Since by hypothesis 
 $f_1 = 1\cdot f_1= p_1\cdot f_1 \sim q_1\cdot g_1 = 1\cdot g_1 =g_1$, then $f_1;\iota_1 \sim g_1;\iota_1$.

If $q_1\not=1$, by hypothesis, we have that $q_2\cdot g_2\sim p_2\cdot f_2=0 \cdot f_2 = \star = q_2\cdot\star$. If $q_2\not=0$, by cancellativity we have, 
\begin{equation}\tag{$\heartsuit$}\label{eq:heart}
g_2\sim\star \text{.}\end{equation}
The following derivation
    \begin{align}
        \langle g_1,g_2\rangle_{(q_1,q_2)}&= \, \diagpX{q_1}{}; (g_1\piu\, (\diagpX{\frac{q_2}{1-q_1}\,}{};(g_2\piu\, \bangp{})) ) \tag{Lemma~\ref{lemma:diagpq3}}\\
        &\sim \, \diagpX{q_1}{}; (g_1\piu\, (\diagpX{\frac{q_2}{1-q_1}\,}{};(\star \piu\, \bangp{})) ) \tag{\eqref{eq:heart}}\\
        &= \, \diagpX{q_1}{}; (g_1\piu\, (\diagpX{\frac{q_2}{1-q_1}\,}{};(\,\bangp{};\cobang{} \piu\, \bangp{})) ) \tag{Def. $\star$}\\
        &=\, \diagpX{q_1}{}; (g_1\piu\, (\,\bangp{};(\,\cobang{}\piu \id{0}) ))\tag{\ref{eq:diagp nat}}\\
        &= \, (\diagpX{q_1}{}; (g_1\piu\, \,\bangp{})); (\id{} \oplus \cobang{} )\tag{SMC} \\
        &= (q_1\cdot g_1) ; \iota_1 \tag{Lemma~\ref{lemma:initialproperties}}
    \end{align}
proves that $\langle g_1,g_2\rangle_{(q_1,q_2)}\sim q_1\cdot g_1;\iota_1$. By means of Lemma~\ref{lemma:diagpq3} and Lemma~\ref{lemma:initialproperties}, one can easily prove that $\langle f_1,f_2\rangle_{(p_1,p_2)}=\iota_1^{Z\oplus Z};(f_1\oplus f_2) =f_1;\iota_1 $. Now, by hypothesis, $f_1= 1\cdot f_1 = p_1\cdot f_1 \sim q_1\cdot g_1$. We can thus conclude that $\langle g_1,g_2\rangle_{(q_1,q_2)}\sim \langle f_1,f_2\rangle_{(p_1,p_2)}$. The remaining cases are proved similarly.
\end{proof}


\begin{comment}
    

\begin{proof}\label{proof-lemma: convex products quoziente}\marginpar{OLD}
The cases where $p_i$ or $q_i$ are equal to $1$ are proved similarly, we show the case where $p_1=1$, hence $p_2=0$. If $q_1=1$, then by Lemma~\ref{lemma:diagpq3} easily follows that $\langle f_1,f_2\rangle_{(1,0)}=f_1;\iota_1$ and $\langle g_1,g_2\rangle_{(1,0)}=g_1;\iota_1$, hence the result follows by hypothesis and the fact that $\sim$ is a congruence.

If $q_1\not=1$, cancellativity of $\sim$ and $q_2\cdot g_2\sim p_2\cdot f_2=\star=q_2\cdot\star$ imply that $g_2\sim\star$. Then,
    \begin{align}
        \langle g_1,g_2\rangle_{(q_1,q_2)}&= \, \diagpX{q_1}{}; (g_1\piu\, (\diagpX{\frac{q_2}{1-q_1}\,}{};(g_2\piu\, \bangp{})) ) \tag{Lemma~\ref{lemma:diagpq}.\ref{lemma:diagpq3}}\\
        &\sim \, \diagpX{q_1}{}; (g_1\piu\, (\diagpX{\frac{q_2}{1-q_1}\,}{};(\star \piu\, \bangp{})) ) \tag{$\sim$ congruence}\\
        &= \, \diagpX{q_1}{}; (g_1\piu\, (\diagpX{\frac{q_2}{1-q_1}\,}{};(\,\bangp{};\cobang{} \piu\, \bangp{})) ) \tag{Def. $\star$}\\
        &=\, \diagpX{q_1}{}; (g_1\piu\, (\,\bangp{};(\,\cobang{}\piu \id{0}) ))\tag{\ref{eq:diagp nat}}\\
        &= \, (\diagpX{q_1}{}; (g_1\piu\, \,\bangp{})); (\id{} \oplus \cobang{} )\tag{SMC} \\
        &= q_1\cdot g_1; \iota_1 \tag{Lemma~\ref{lemma:initialproperties}.\ref{lemma:iotai1}}
    \end{align}
    implies $\langle g_1,g_2\rangle_{(q_1,q_2)}\sim q_1\cdot g_1;\iota_1$. Since by Lemma~\ref{lemma:diagpq}.\ref{lemma:diagpq3} and Lemma~\ref{lemma:initialproperties}.\ref{lemma:iotai1}, $\langle f_1,f_2\rangle_{(p_1,p_2)}=\iota_1^{Z\oplus Z};(f_1\oplus f_2) =f_1;\iota_1 $ and, by hypothesis, $f_1\sim q_1\cdot g_1$, the post-composition with $\iota_1$ proves the result. The cases where $p_i$ or $q_i$ are equal to $0$ are proved similarly.
   
   
    Now consider the cases where $p_i,q_i\in(0,1)$. Define $r_i\defeq \begin{cases}f_i &\text{if }p_i\leq q_i\\ g_i & \text{else} \end{cases}$, $u_i\defeq \begin{cases}p_i &\text{if }p_i\leq q_i\\ q_i & \text{else} \end{cases}$ and consider $\langle r_1,r_2\rangle_{(u_1,u_2)}$. We prove that $\langle f_1,f_2\rangle_{(p_1,p_2)}\sim\langle r_1,r_2\rangle_{(u_1,u_2)}\sim\langle g_1,g_2\rangle_{(q_1,q_2)}$. Without loss of generality, we can assume that $p_1\leq q_1$  and $q_2\leq p_2$, hence $\langle r_1,r_2\rangle_{(u_1,u_2)}=\langle f_1,g_2\rangle_{(p_1,q_2)}$.  The first equality is obtained as follows:
  \begin{align*}
        \langle f_1,f_2\rangle_{(p_1,p_2)}&=\, \diagpX{p_1}{}; (f_1\piu\, (\diagpX{\frac{p_2}{1-p_1}\,}{};(f_2\piu\, \bangp{})) ) \tag{Lemma~\ref{lemma:diagpq}.\ref{lemma:diagpq3}}\\
        &\sim \, \diagpX{p_1}{}; (f_1\piu\, (\diagpX{\frac{q_2}{1-p_1}\,}{};(g_2\piu\, \bangp{})) ) \tag{$\ast$}\\
        &=\, \langle f_1,g_2\rangle_{(p_1,q_2)} \tag{Lemma~\ref{lemma:diagpq}.\ref{lemma:diagpq3}}
  \end{align*}
 where $(\ast)$ follows from cancellativity of $\sim$ and the fact that $p_2\cdot f_2\sim q_2\cdot g_2$  and that $p_2\cdot f_2= (1-p_1)\cdot (\frac{p_2}{1-p_1}\cdot f_2)$ while $q_2\cdot g_2= (1-p_1)\cdot (\frac{q_2}{1-p_1}\cdot g_2)$.
Hence $\langle f_1,f_2\rangle_{(p_1,p_2)}\sim\langle f_1,g_2\rangle_{(p_1,q_2)}$. The second equality is obtained as follows:
\begin{align*}
    \langle g_1,g_2\rangle_{(q_1,q_2)}&=\, \diagpX{q_1}{}; (g_1\piu\, (\diagpX{\frac{q_2}{1-q_1}\,}{};(g_2\piu\, \bangp{})) ) \tag{Lemma~\ref{lemma:diagpq}.\ref{lemma:diagpq3}}\\
    &= \, \diagpX{q_2}{};(g_2\piu\, (\diagpX{\frac{q_1}{1-q_2}\,}{};(g_1\piu\, \bangp{})) ) \tag{\ref{eq:diagp assoc} + \ref{eq:diagp symmetry}}\\
    &\sim \, \diagpX{q_2}{};(g_2\piu\, (\diagpX{\frac{p_1}{1-q_2}\,}{};(f_1\piu\, \bangp{})) ) \tag{$\ast$}\\
    &=\,  \diagpX{p_1}{}; (f_1\piu\, (\diagpX{\frac{q_2}{1-p_1}\,}{};(g_2\piu\, \bangp{})) ) \tag{\ref{eq:diagp assoc} + \ref{eq:diagp symmetry}}     \\
    &=\, \langle f_1,g_2\rangle_{(p_1,q_2)} \tag{Lemma~\ref{lemma:diagpq}.\ref{lemma:diagpq3}}
\end{align*}
where $(\ast)$ is obtained as above. Hence $\langle g_1,g_2\rangle_{(q_1,q_2)}\sim\langle f_1,g_2\rangle_{(p_1,q_2)}$.
\end{proof}

\end{comment}







\begin{proof}[Proof of Proposition~\ref{cor: quotient category is convex biproduct }]
   By construction, every object in $\CatTapeC$ carries a commutative and coherent monoid and co-pca structure. Since $\sim$ is a congruence relation, the equational axioms in Table~\ref{fig:freestrictfccat} and \ref{fig:freecopcacat}  hold in $\CatTapeC_\sim$. Thus, thanks to Proposition~\ref{prop:A}, to prove that $\CatTapeC_\sim$ is a convex biproduct category, it is enough to prove that it has jointly monic projections.
  
  
First, we consider two arrows of the form $[\t]_\sim,[\s]_\sim\colon U\to Q\oplus R$ where $U$ is an object of $\Cat{C}$. We assume that $[\t]_\sim;[\pi_1]_\sim = [\s]_\sim;[\pi_1]_\sim$ and $[\t]_\sim;[\pi_2]_\sim = [\s]_\sim;[\pi_2]_\sim$. By Lemma~\ref{decomposition}, $\t=\,\diaggen{p,q}{};(\t_1\oplus \t_2)$ and $\s=\,\diaggen{p',q'}{};(\s_1\oplus \s_2)$. By the assumption and Lemma~\ref{lemma:diagpq}, it holds that $p\cdot \t_1 \sim p'\cdot \s_1$ and $q\cdot \t_2 \sim q'\cdot \s_2$. Hence, by Lemma~\ref{lemma: convex products quoziente}, $\langle \t_1,\t_2\rangle_{(p,q)}\sim\langle \s_1,\s_2\rangle_{(p',q')}$ and thus $[\t]_\sim=[\s]_\sim$.
    
 The general case $[\t]_\sim,[\s]_\sim\colon P\to Q\oplus R$ where $P=\bigoplus_{i=1}^n U_i$ is proved similarly by relying on the fact that $\CatTapeC_\sim$  also has coproducts thanks to Fox's Theorem. Indeed, by the universal property of coproducts, we can write $\t = [\t_1,\dots,\t_n]$ and $\s = [\s_1,\dots,\s_n]$ for some $\t_i,\s_i\colon U_i \to Q\oplus R$. Now, from $[\t]_\sim;[\pi_1]_\sim = [\s]_\sim;[\pi_1]_\sim$ and $[\t]_\sim;[\pi_2]_\sim = [\s]_\sim;[\pi_2]_\sim$ it follows that $[\t_i]_\sim;[\pi_1]_\sim = [\s_i]_\sim;[\pi_1]_\sim$ and $[\t_i]_\sim;[\pi_2]_\sim = [\s_i]_\sim;[\pi_2]_\sim$ for all $i=1,\dots,n$. By the previous case, we have that $[\t_i]_\sim=[\s_i]_\sim$ for all $i=1,\dots,n$ and hence $[\t]_\sim=[\s]_\sim$, given that $[[\t_1,\dots,\t_n]]_\sim=[[\t_1]_\sim,\dots,[\t_n]_\sim]$ and $[[\s_1,\dots,\s_n]]_\sim=[[\s_1]_\sim,\dots,[\s_n]_\sim]$.

 The functor $Q_\sim\colon \CatTapeC\to \CatTapeC_\sim$ is a morphism of convex biproduct categories since it preserves finite coproducts.
\end{proof}

\begin{comment}
\marginpar{Da qui in poi possiamo oscurare.}

\cipr{Problemi trovati:
\begin{itemize}
    \item Il Lemma \ref{prop: convex products quoziente} non funziona: dicevamo che se $p_2\cdot f_2\sim q_2\cdot g_2$ allora $\frac{p_2}{1-p_1}\cdot f_2\sim \frac{q_2}{1-p_1}\cdot g_2$ ma questo non sembra essere derivabile. Quello che sappiamo è che se $p_2\cdot f_2\sim q_2\cdot g_2$ allora $p\cdot p_2\cdot f_2\sim p\cdot q_2\cdot g_2$ per ogni $p\in [0,1]$ ma questo non è sufficiente perché $\frac{1}{1-p_1}$ è maggiore di 1. 
    \item Diversamente, se $p_2\cdot f_2 = q_2\cdot g_2$ allora $\frac{p_2}{1-p_1}\cdot f_2 = \frac{q_2}{1-p_1}\cdot g_2$. Questo lo si ottiene per cancellatività: visto che $p_2\cdot f_2 = (1-p_1)\cdot \frac{p_2}{1-p_1}\cdot f_2$ e $q_2\cdot g_2 = (1-p_1)\cdot \frac{q_2}{1-p_1}\cdot g_2$, dunque la cancelatività (con $1-p_1$) implica che $\frac{p_2}{1-p_1}\cdot f_2 = \frac{q_2}{1-p_1}\cdot g_2$.
    \item Quindi quello che ci farebbe tornare il discorso è la cancellatività per $\sim$: se $p\cdot f \sim p\cdot g$ allora $f\sim g$ per ogni $p\in (0,1)$. Se questa proprietà è soddisfatta allora il Lemma \ref{prop: convex products quoziente} è vero e dunque anche la Proposition \ref{cor: quotient category is convex biproduct } è vera.
    \item Cercare di dimostrare la cancellatività per la $\sim$ generata dall'assioma \ref{eq:natmultiplexer}, ci siamo accorti che questo assioma non vale per i tape "parziali". Infatti, se prendiamo $p\cdot \id{A}$ otterremmo qualcosa del tipo $p\cdot \id{A} = p^2\cdot \id{A}$.
    \item qui sotto riporto le vechie dim lemma e tma su $\sim$ Sopra ho messo le versioni nuove aggiungendo cancellativa.
\end{itemize}
}



\begin{proposition}\label{prop: convex products quoziente}
    Given a category $\Cat{C}$ and a congruence relation $\sim$ on $\CatTapeC$, then the functor $Q_\sim\colon\CatTapeC\to \CatTapeC_\sim$ satisfies the following property: if $[p_1\cdot f_1]_\sim =[q_1\cdot g_1]_\sim$ and $[p_2\cdot f_2]_\sim =[q_2\cdot g_2]_\sim$  then $[\langle f_1,f_2\rangle_{(p_1,p_2)}]_\sim= [\langle g_1,g_2\rangle_{(q_1,q_2)}]_\sim$.
\end{proposition}
\begin{proof}\label{proof-prop: convex products quoziente}
    Before proceeding with the proof, we make the trivial observation that if $[f]_\sim=[g]_\sim$ then $[p\cdot f]_\sim=[p\cdot g]_\sim$ for all $p\in [0,1]$. Indeed,
\begin{align}
    p\cdot f&=\, \diagpX{p}{}; (f\piu\, \bangp{}) \tag{Thm. \ref{thm:TCconvexbiproductcategory}}\\
    &\sim \, \diagpX{p}{}; (g\piu\, \bangp{}) \tag{by $[f]_\sim=[g]_\sim$}\\
    &=\, p\cdot g \notag
\end{align}
hence $[p\cdot f]_\sim=[p\cdot g]_\sim$.
The cases with $p=0$ and $p=1$ are trivial.
  Now, define  $r_i\defeq \begin{cases}f_i &\text{if }p_i\leq q_i\\ g_i & \text{else} \end{cases}$, $u_i\defeq \begin{cases}p_i &\text{if }p_i\leq q_i\\ q_i & \text{else} \end{cases}$ and consider $\langle r_1,r_2\rangle_{(u_1,u_2)}$. We prove that $[\langle f_1,f_2\rangle_{(p_1,p_2)}]_\sim=[\langle r_1,r_2\rangle_{(u_1,u_2)}]_\sim=[\langle g_1,g_2\rangle_{(q_1,q_2)}]_\sim$. Without loss of generality, we can assume that $p_1\leq q_1$  and $q_2\leq p_2$, hence $\langle r_1,r_2\rangle_{(u_1,u_2)}=\langle f_1,g_2\rangle_{(p_1,q_2)}$.  The first equality is obtained as follows:
  \begin{align*}
        \langle f_1,f_2\rangle_{(p_1,p_2)}&=\, \diagpX{p_1}{}; (f_1\piu\, (\diagpX{\frac{p_2}{1-p_1}\,}{};(f_2\piu\, \bangp{})) ) \tag{Thm. \ref{thm:TCconvexbiproductcategory}}\\
        &\sim \, \diagpX{p_1}{}; (f_1\piu\, (\diagpX{\frac{q_2}{1-p_1}\,}{};(g_2\piu\, \bangp{})) ) \tag{by $[p_2\cdot f_2]_\sim =[q_2\cdot g_2]_\sim$}\\
        &=\, \langle f_1,g_2\rangle_{(p_1,q_2)}
  \end{align*}

hence $[\langle f_1,f_2\rangle_{(p_1,p_2)}]_\sim=[\langle f_1,g_2\rangle_{(p_1,q_2)}]_\sim$. The second equality is obtained as follows:
\begin{align*}
    \langle g_1,g_2\rangle_{(q_1,q_2)}&=\, \diagpX{q_1}{}; (g_1\piu\, (\diagpX{\frac{q_2}{1-q_1}\,}{};(g_2\piu\, \bangp{})) ) \tag{Thm. \ref{thm:TCconvexbiproductcategory}}\\
    &= \, \diagpX{q_2}{};(g_2\piu\, (\diagpX{\frac{q_1}{1-q_2}\,}{};(g_1\piu\, \bangp{})) ) \tag{\ref{eq:diagp assoc} + \ref{eq:diagp symmetry}}\\
    &\sim \, \diagpX{q_2}{};(g_2\piu\, (\diagpX{\frac{p_1}{1-q_2}\,}{};(f_1\piu\, \bangp{})) ) \tag{by $[p_1\cdot f_1]_\sim =[q_1\cdot g_1]_\sim$}\\
    &=  \diagpX{p_1}{}; (f_1\piu\, (\diagpX{\frac{q_2}{1-p_1}\,}{};(g_2\piu\, \bangp{})) ) \tag{\ref{eq:diagp assoc} + \ref{eq:diagp symmetry}}     \\
    &=\, \langle f_1,g_2\rangle_{(p_1,q_2)}
\end{align*}
hence $[\langle g_1,g_2\rangle_{(q_1,q_2)}]_\sim=[\langle f_1,g_2\rangle_{(p_1,q_2)}]_\sim$.
\end{proof}

\begin{proof}[Proof of Proposition \ref{cor: quotient category is convex biproduct }]\label{proof-cor: quotient category is convex biproduct }
    Since $\sim$ is a congruence relation, $\CatTapeC_\sim$ is a monoidal category and $Q_\sim$ is a monoidal functor. The object $0$ is initial and terminal in $\CatTapeC_\sim$ since it is so in $\CatTapeC$. The diagram  $A\overset{[\iota_1^{A\piu B}]_\sim}{\rightarrow}A\piu B\overset{[\iota_2^{A\piu B}]_\sim}{\leftarrow} B$ is the coproduct of $A$ and $B$ since for every pair of arrows $[f]_\sim\colon A\to X$ and $[g]_\sim:B\to X$ there exists the arrow $[(f\piu g);\codiag{X}]_\sim$ that gives $[\iota_1^{A\piu B}]_\sim;[(f\piu g);\codiag{X}]_\sim=[f]_\sim$ and $[\iota_B]_\sim;[(f\piu g);\codiag{X}]_\sim=[g]_\sim$. Such arrow is the unique satisfying this property, indeed if $[h]\colon A\piu B\to X$ satisfies $[\iota_1^{A\piu B}]_\sim;[h]_\sim=[f]_\sim$ and $[\iota_2^{A\piu B}]_\sim;[h]_\sim=[g]_\sim$, then coproducts in $\CatTapeC$ imply that there exist $f',g'$ such that $h\coloneqq [f',g'];\codiag{X}$ and then
    \begin{align*}
[\iota_1^{A\piu B}]_\sim;[h]_\sim&=[\iota_1^{A\piu B}]_\sim;[[f',g'];\codiag{X}]_\sim \tag{$h\coloneqq [f',g'];\codiag{X}$}\\
&=[f']_\sim \tag{$\sim$ congruence relation}\\
&=[f]_\sim \tag{by hp.}
        \end{align*}
        similarly $[\iota_2^{A\piu B}]_\sim;[h]_\sim=[g']_\sim=[g]_\sim$. Hence, we obtain
        \begin{align*}
            [h]_\sim&=[[f',g'];\codiag{X}]_\sim \tag*{}\\
            &=[f']_\sim\piu [g']_\sim;[\codiag{X}]_\sim \tag{$\sim$ congruence relation}\\
            &=[f]_\sim\piu [g]_\sim;[\codiag{X}]_\sim \tag{above equalities}\\
            &=[(f\piu g);\codiag{X}]_\sim \tag{$\sim$ congruence relation}
        \end{align*}
    $\CatTapeC_\sim$ is also enriched over convex algebras. Indeed, for all $X,Y\in \CatTapeC_\sim$, one defines $[f]_\sim +_p [g]_\sim \defeq [f +_p g]_\sim$ for all $[f]_\sim,[g]_\sim \in \CatTapeC_\sim(X,Y)$ and $p\in [0,1]$. This is well defined since if $[f]_\sim=[f']_\sim$ and $[g]_\sim=[g']_\sim$ then
    \begin{align*}
        [f +_p g]_\sim &= [\diagpX{p}{X}; (f\piu g);\codiag{Y}]_\sim \tag{Thm.~\ref{thm:TCconvexbiproductcategory}}\\
        &=[\diagpX{p}{X}]_\sim;( [f]_\sim\piu [g]_\sim);[\codiag{Y}]_\sim \tag{$\sim$ congruence relation}\\
        &=[\diagpX{p}{X}]_\sim;( [f']_\sim\piu [g']_\sim);[\codiag{Y}]_\sim \tag{by $[f]_\sim=[f']_\sim$ and $[g]_\sim=[g']_\sim$}\\
        &=[f' +_p g']_\sim
        \end{align*}
     
While $\star_{X,Y}$ is just $[\star_{X,Y}]_\sim$. Axioms of PCA are satisfied since they are so in $\CatTapeC$ and $\sim$ is a congruence relation. Now we prove that $A\overset{[\pi_1]_\sim}{\leftarrow}A\piu B\overset{[\pi_2]_\sim}{\rightarrow} B$ is a binary convex product also in $\CatTapeC$. Indeed, given $[f]_\sim \colon X\to A$ and $[g]_\sim \colon X\to B$, $p_1,p_2\in [0,1]$ such that $p_1+p_2\leq 1$, we prove that there is a unique $[h]_\sim \colon X\to A\oplus B$ such that $[h]_\sim;[\pi_1]_\sim = p_1\cdot [f]_\sim$ and $[h]_\sim;[\pi_2]_\sim = p_2\cdot [g]_\sim$. Existence is provided by $h\defeq \langle f,g\rangle_{(p_1,p_2)}$, for uniqueness, suppose that there is also $[h']_\sim$ with the same property. Then, by Theorem~\ref{thm:TCconvexbiproductcategory}, $h'=\langle f',g'\rangle_{(q_1,q_2)}$ for some $f'\colon X\to A$, $g'\colon X\to B$ and $q_1,q_2\in [0,1]$ such that $q_1+q_2\leq 1$. Now, since $[h']_\sim;[\pi_1]_\sim = [q_1\cdot f']_\sim=[p_1\cdot f]_\sim$ and $[h']_\sim;[\pi_2]_\sim = [q_2\cdot g' ]_\sim = p_2\cdot [g]_\sim$, by Proposition~\ref{prop: convex products quoziente} we have that $[h']_\sim=[\langle f,g\rangle_{(p_1,p_2)}]_\sim$. Hence, $\langle [f],[g]\rangle_{(p_1,p_2)}\coloneqq[\langle f,g\rangle_{(p_1,p_2)}]$ is the unique arrow with the required property. Observe also that, by definition, $[f]_\sim +_p [g]_\sim = [\diagp{}]; ([f]_\sim \piu [g]_\sim); [\codiag{}]_\sim$. The fact that $Q_\sim$ is a convex biproduct functor because is PCA-enriched and preserves coproducts.
 \end{proof}













\begin{proposition}%\label{cor: quotient category is convex biproduct2}
Let $\Cat{C}$ be a category and $\sim$ be a congruence relation (w.r.t.\ $;$ and $\oplus$) on $\CatTapeC$ which is cancellative: for every $p\in (0,1)$, $p\cdot f\sim p\cdot g$ then $f\sim g$. Let $\CatTapeC_\sim$ be the category obtained as the quotient of $\CatTapeC$ by $\sim$ and $Q_\sim \colon \CatTapeC\to \CatTapeC_\sim$ be the functor mapping each arrow to its $\sim$-equivalence class.
Then $\CatTapeC_\sim$ is a convex biproduct category and $Q_\sim$ is a morphism of convex biproduct categories.
\end{proposition}

%%%%%%VECCHIE PROVE DI DIMOSTRAZIONE TC_\sim CBC CHE NON USANO J MONO PROJ E LA NUOVA VERSIONE DEL LEMMA DI DEOMPOSIZIONE, SI POSSONO PURE CANCELLARE%%%%%%%%%%%

    

\begin{proof}[Proof of Proposition \ref{cor: quotient category is convex biproduct2}]
    By construction, every object in $\CatTapeC$ is a commutative and coherent monoid and co-pca object. Since $\sim$ is a congruence relation, the equational axioms in Table~\ref{fig:freestrictfccat} and \ref{fig:freecopcacat}  hold in $\CatTapeC_\sim$. Hence, by Fox's Theorem and Lemma~\ref{lemma:enrichmentpca}, $\CatTapeC_\sim$ is a pca-enriched category with coproducts and zero object.    
    We now prove that for every $X_1,X_2$ in $\CatTapeC_\sim$, the object $X_1\oplus X_2$ with the projections $\pi_1= [\id{X_1}\oplus\, \bangp{X_2}]_\sim$ and $\pi_2=[\,\bangp{X_1}\oplus \id{X_2}]_\sim$ is a convex product. Consider $[f]_\sim\colon U\to X_1$, $[g]_\sim\colon U\to X_2$ and $p_1,p_2\in [0,1]$ such that $p_1+p_2\le 1$, where $U\in \Cat{C}$. The existence property of convex products follows considering $h \defeq [\langle f, g\rangle_{(p_1,p_2)}]_\sim$. Indeed, a simple verification shows that $h;\pi_1= [p_1\cdot f]_\sim= p_1\cdot[f]_\sim$ and $h;\pi_2= [p_2\cdot g]_\sim= p_2\cdot[g]_\sim$. Assume now that there exists another $h'$ such that $h';\pi_1=p_1\cdot [f]$ and $h'\cdot \pi_2=p_2\cdot[g]$. By Lemma~\ref{lemma:division}, $h'$ must be the equivalence class of an arrow in one of the three forms indicated. Consider the case $h'=[\diagq{};f'\oplus g']_\sim=[\langle f',g'\rangle_{q,1-q}]$ and observe that $h';\pi_1= [q\cdot f']_\sim=[p_1\cdot f]_\sim $ and $h';\pi_2=[(1-q)\cdot g']_\sim=[p_2\cdot g]_\sim$. Hence, Lemma~\ref{lemma: convex products quoziente} implies $h'=h$. Now assume $h'$ is the equivalence class of one of the other forms provided by Lemma~\ref{lemma:division}. We prove the result for the case $h'=[f'\oplus \cobang{}]_\sim$, the other cases are proved similarly. 
    Since $h';\pi_1= [f']_\sim= [p_1\cdot f]_\sim$  and $h';\pi_2= [\star]= [p_2\cdot g]_\sim$, cancellativity of $\sim$ implies that $g\sim \star$. Then,
    \begin{align}
        \langle f,g\rangle_{(p_1,p_2)}&=\, \diagpX{p_1}{}; (f\piu\, (\diagpX{\frac{p_2}{1-p_1}\,}{};(g\piu\, \bangp{})) ) \tag{Thm. \ref{thm:TCconvexbiproductcategory}}\\
        &\sim \,\diagpX{p_1}{}; (f\piu\, (\diagpX{\frac{p_2}{1-p_1}\,}{};(\star\piu\, \bangp{})) ) \notag\\
        &= p_1\cdot f; \iota_1 \notag\\
        &\sim f'; \iota_1= f'\oplus \cobang \notag
    \end{align}
    implies that $h'=h$. For the general case $f\colon P\to X_1$ and $g\colon P\to X_2$ where $P=\bigoplus_{i=1}^n U_i$, we rely on the fact that $\CatTapeC_\sim$ has coproducts. Indeed, $h$ is the equivalence class of the copairing of the arrows obtained from the convex product with arrows $f_i\defeq \iota_i;f$ and $g_i\defeq \iota_i;g$ for $i=1,\dots,n$ and costants $p_1,p_2$. If there exists another $h'$ such that $h';\pi_1=p_1\cdot [f]$ and $h'\cdot \pi_2=p_2\cdot[g]$, then for every $i$ it holds by the previous case that $\iota_i;h'=\iota_i;h$. Finally, the universal property of coproducts imply $h=h'$. $Q_\sim$ is a convex biproduct functor because is PCA-enriched and preserves coproducts.
\end{proof} 
{\color{red} Propongo questa modifica nel tma \ref{thm:TCconvexbiproductcategory} e nel lemma di decomposizione: sostituire $U_i$ con $R_i$, che resta vera e si può usare meglio nelle due dimostrazioni di sopra. Con questa assunzione l'ultima dimostrazione si semplifica come segue:}


\begin{proof}[New Proof of Proposition \ref{cor: quotient category is convex biproduct }]
    By construction, every object in $\CatTapeC$ is a commutative and coherent monoid and co-pca object. Since $\sim$ is a congruence relation, the equational axioms in Table~\ref{fig:freestrictfccat} and \ref{fig:freecopcacat}  hold in $\CatTapeC_\sim$. Hence, by Fox's Theorem and Lemma~\ref{lemma:enrichmentpca}, $\CatTapeC_\sim$ is a pca-enriched category with coproducts and zero object.    
    We now prove that for every $X_1,X_2$ in $\CatTapeC_\sim$, the object $X_1\oplus X_2$ with the projections $\pi_1= [\id{X_1}\oplus\, \bangp{X_2}]_\sim$ and $\pi_2=[\,\bangp{X_1}\oplus \id{X_2}]_\sim$ is a convex product. Consider $[f]_\sim\colon U\to X_1$, $[g]_\sim\colon U\to X_2$ and $p_1,p_2\in [0,1]$ such that $p_1+p_2\le 1$, where $U\in \Cat{C}$. The existence property of convex products follows considering $h \defeq [\langle f, g\rangle_{(p_1,p_2)}]_\sim$. Indeed, a simple verification shows that $h;\pi_1= [p_1\cdot f]_\sim= p_1\cdot[f]_\sim$ and $h;\pi_2= [p_2\cdot g]_\sim= p_2\cdot[g]_\sim$. Assume now that there exists another $h'$ such that $h';\pi_1=p_1\cdot [f]$ and $h'\cdot \pi_2=p_2\cdot[g]$. By Lemma~\ref{decomposition}, $h'$ must be the equivalence class of an arrow of the form $\diagpn{\vec{q}\;\;}{}{2}; {f'\oplus g'}$, with $\vec{q}=(q_1,q_2)$. Hence, $h';\pi_1= [q_1\cdot f']_\sim=[p_1\cdot f]_\sim $ and $h';\pi_2=[q_2\cdot g']_\sim=[p_2\cdot g]_\sim$, and Lemma~\ref{lemma: convex products quoziente} implies $h'=h$. For the general case $f\colon P\to X_1$ and $g\colon P\to X_2$ where $P=\bigoplus_{i=1}^n U_i$, we rely on the fact that $\CatTapeC_\sim$ has coproducts. Indeed, $h$ is the equivalence class of the copairing of the arrows obtained from the convex product with arrows $f_i\defeq \iota_i;f$ and $g_i\defeq \iota_i;g$ for $i=1,\dots,n$ and costants $p_1,p_2$. If there exists another $h'$ such that $h';\pi_1=p_1\cdot [f]$ and $h'\cdot \pi_2=p_2\cdot[g]$, then for every $i$ it holds by the previous case that $\iota_i;h'=\iota_i;h$. Finally, the universal property of coproducts imply $h=h'$. $Q_\sim$ is a convex biproduct functor because is PCA-enriched and preserves coproducts.
\end{proof} 

\end{comment}
